
\section{Predictable Quadratic Variation Concentration}

We control the predictable quadratic variation of \(M_n^{\text{RR}}\) by applying
\eqref{eq:external-markov-conc} to a centered quadratic functional of the chain.
% This is the RR analogue of Lemmas 22--23 of Samsonov et al. (2025).

\textbf{Predictable quadratic variation.} The increments
\begin{equation}
\begin{aligned}
\Delta M_l^{\text{RR}} = \mathcal{Q}_l^{\text{RR}} \, (\widehat{\varepsilon}(Z_l) - \mathsf{Q} \widehat{\varepsilon}(Z_{l - 1}))
\end{aligned}
\end{equation}
are \(\mathcal{F}_l\)-martingale differences. The Markov property gives
\begin{equation}
\begin{aligned}
\mathbb{E}[\widehat{\varepsilon}(Z_l) | \mathcal{F}_{l - 1}] = \mathsf{Q} \widehat{\varepsilon}(Z_{l - 1})
\end{aligned}
\end{equation}
% Direct computation:
\begin{equation}
\begin{aligned}
\mathbb{E}[\Delta M_l^{\text{RR}} \, (\Delta M_l^{\text{RR}})^\top \, | \, \mathcal{F}_{l - 1}]
% &= \mathcal{Q}_l^{\text{RR}}
%    \mathbb{E}[\widehat{\varepsilon}(Z_l)\widehat{\varepsilon}(Z_l)^\top
%      \mid\mathcal{F}_{l-1}](\mathcal{Q}_l^{\text{RR}})^\top \\
% &\quad-\mathcal{Q}_l^{\text{RR}}\mathsf{Q}\widehat{\varepsilon}(Z_{l-1})
%    (\mathsf{Q}\widehat{\varepsilon})(Z_{l-1})^\top
%    (\mathcal{Q}_l^{\text{RR}})^\top \\
  &= \mathcal{Q}_l^{\text{RR}} \, \mathcal{V}_{\varepsilon}(Z_{l - 1}) \, (\mathcal{Q}_l^{\text{RR}})^\top,
\end{aligned}
\end{equation}
where
\begin{equation}
\begin{aligned}
\mathcal{V}_{\varepsilon}(z)
  \coloneqq \mathsf{Q}(\widehat{\varepsilon} \widehat{\varepsilon}^\top)(z)
   - (\mathsf{Q} \widehat{\varepsilon})(z) \, (\mathsf{Q} \widehat{\varepsilon})(z)^\top.
\end{aligned}
\label{eq:bar-eps-def}
\end{equation}
% The cross-terms cancel by the Markov property. This is a matrix-valued
% conditional covariance of the Poisson martingale increment; the vector noise
% remains denoted by \(\varepsilon\).

\begin{lemma}\label{lem:poisson-covariance-identity}
  \textbf{(Poisson covariance identity and absolute convergence.)}
  Assume \textbf{UGE 1}, \(\pi(\varepsilon) = 0\), and
  \(\lVert \varepsilon \rVert_\infty < \infty\). Then \(\Sigma_{\varepsilon}^{(\text{M})}\) is absolutely convergent in operator
  norm, and
\begin{equation}
\begin{aligned}
\pi(\mathcal{V}_{\varepsilon}) = \Sigma_{\varepsilon}^{(\text{M})} =: \Sigma.
\end{aligned}
\label{eq:poisson-cov-identity}
\end{equation}
\end{lemma}

\textit{Proof.} Boundedness and UGE give absolute convergence. For \(j \ge 1\),
\begin{equation}
\begin{aligned}
\lVert \mathbb{E}_\pi [\varepsilon(Z_0) \varepsilon(Z_j)^\top] \rVert
  = \lVert \pi(\varepsilon \, (\mathsf{Q}^j \varepsilon)^\top) \rVert
  \le \lVert \varepsilon \rVert_\infty \lVert \mathsf{Q}^j \varepsilon \rVert_\infty
  \le 2 \lVert \varepsilon \rVert_\infty^2 (1 / 4)^{\left\lfloor j / t_{\text{mix}} \right\rfloor},
\end{aligned}
\label{eq:SigmaM-abs-one-sided}
\end{equation}
and similarly for the transposed covariance. Set \(h \coloneqq \widehat{\varepsilon}\) and \(r \coloneqq \mathsf{Q} h = h - \varepsilon\). Since
\(\pi \mathsf{Q} = \pi\),
\begin{equation}
\begin{aligned}
\pi(\mathcal{V}_{\varepsilon})
  &= \pi(h h^\top) - \pi(r r^\top)
  &= \pi(\varepsilon \varepsilon^\top)
     + \pi(\varepsilon r^\top) + \pi(r \varepsilon^\top).
\end{aligned}
\end{equation}
The Poisson series gives \(r = \sum_{j \ge 1} \mathsf{Q}^j \varepsilon\) in sup-norm:
\begin{equation}
\begin{aligned}
\pi(\varepsilon r^\top)
  &= \sum_{j \ge 1}
     \mathbb{E}_\pi [\varepsilon(Z_0) \varepsilon(Z_j)^\top], \\
\pi(r \varepsilon^\top)
  &= \sum_{j \ge 1}
     \mathbb{E}_\pi [\varepsilon(Z_j) \varepsilon(Z_0)^\top].
\end{aligned}
\end{equation}
Substitution proves \eqref{eq:poisson-cov-identity}. \(\square\)
% This is the Poisson-solution covariance identity used by Samsonov et al.
% (2025, Eq. (10)) and by the Markov-chain CLT formulation of
% Douc--Moulines--Priouret--Soulier (2018, Theorem 21.2.5).

Thus
\begin{equation}
\begin{aligned}
\langle M^{\text{RR}} \rangle_n
  &\coloneqq \sum_{l = 2}^{n - 1}
      \mathbb{E}[\Delta M_l^{\text{RR}} \, (\Delta M_l^{\text{RR}})^\top | \mathcal{F}_{l - 1}]
  &= \sum_{l = 2}^{n - 1}
      \mathcal{Q}_l^{\text{RR}} \, \mathcal{V}_{\varepsilon}(Z_{l - 1})
      \, (\mathcal{Q}_l^{\text{RR}})^\top.
\end{aligned}
\label{eq:M-RR-bracket}
\end{equation}
By Lemma~\ref{lem:poisson-covariance-identity}, \(\pi(\mathcal{V}_{\varepsilon}) = \Sigma\).

\textbf{Sup-norm bound on \(\mathcal{V}_{\varepsilon}\).} Since
\begin{equation}
\begin{aligned}
\lVert  \widehat{\varepsilon}  \rVert_\infty \le 3 \, t_{\text{mix}} \, \lVert  \varepsilon  \rVert_\infty
\end{aligned}
\end{equation}
and \(\mathsf{Q}\) is a Markov kernel,
\begin{equation}
\begin{aligned}
\lVert  \mathcal{V}_{\varepsilon}  \rVert_\infty
  &\le \lVert  \mathsf{Q}(\widehat{\varepsilon} \widehat{\varepsilon}^\top)  \rVert_\infty
   + \lVert  (\mathsf{Q} \widehat{\varepsilon})(\mathsf{Q} \widehat{\varepsilon})^\top  \rVert_\infty \\
  &\le 2 \, \lVert  \widehat{\varepsilon}  \rVert_\infty^2
   \le 18 \, t_{\text{mix}}^2 \, \lVert  \varepsilon  \rVert_\infty^2.
\end{aligned}
\label{eq:bar-eps-sup}
\end{equation}

\begin{lemma}\label{lem:M-RR-bracket-conc}
  \textbf{(Stationary predictable-variation concentration.)}
  Assume \textbf{UGE 1} and \(\pi(\varepsilon) = 0\), with \(\lVert  \varepsilon  \rVert_\infty < \infty\).
  Let \(C_{\mathcal{Q}}\) be the uniform bound on \(\lVert  \mathcal{Q}_l^{\text{RR}}  \rVert\) from the
  previous lemma. There exists a universal constant \(C_4 > 0\) such that, for
  every \(u \in \mathbb{R}^d\), every \(p \ge 2\), every initial distribution \(\xi\), and every
  \(n \ge 2\),
\begin{equation}
\begin{aligned}
  \mathbb{E}_\xi^{1 / p} \left[
    | u^\top \langle M^{\text{RR}} \rangle_n u - n \, \sigma_n^{2, \text{RR}}(u) |^p
  \right]
    \le C_4 \, C_{\mathcal{Q}}^2 \, \lVert  u  \rVert^2 \, \lVert  \varepsilon  \rVert_\infty^2
       \, t_{\text{mix}}^{5 / 2} \, \sqrt{p \, n}.
\end{aligned}
\label{eq:M-RR-conc}
\end{equation}
\end{lemma}

\textit{Proof.} Write \(h_l(z) \coloneqq u^\top \mathcal{Q}_l^{\text{RR}} \, \mathcal{V}_{\varepsilon}(z) \, (\mathcal{Q}_l^{\text{RR}})^\top u\)
and \(g_l(z) \coloneqq h_l(z) - \pi(h_l)\). Then
\begin{equation}
\begin{aligned}
|h_l(z)|
  \le C_{\mathcal{Q}}^2 \, \lVert  u  \rVert^2 \, \lVert  \mathcal{V}_{\varepsilon}  \rVert_\infty
  \le 18 \, C_{\mathcal{Q}}^2 \, \lVert  u  \rVert^2 \, t_{\text{mix}}^2 \, \lVert  \varepsilon  \rVert_\infty^2,
\quad
\lVert  g_l  \rVert_\infty \le 2 \, |h_l(z)|
  \le 36 \, C_{\mathcal{Q}}^2 \, \lVert  u  \rVert^2 \, t_{\text{mix}}^2 \, \lVert  \varepsilon  \rVert_\infty^2.
\end{aligned}
\end{equation}
By \eqref{eq:M-RR-bracket},
\(n \, \sigma_n^{2, \text{RR}}(u) = \sum_{l = 2}^{n - 1} \pi(h_l)\),
\begin{equation}
\begin{aligned}
u^\top \langle M^{\text{RR}} \rangle_n u - n \, \sigma_n^{2, \text{RR}}(u)
  = \sum_{l = 2}^{n - 1} g_l(Z_{l - 1}).
\end{aligned}
\label{eq:M-RR-conc-decomp}
\end{equation}

Set \(\widetilde{g}_i \coloneqq g_{i + 1}\). Then \(\pi(\widetilde{g}_i) = 0\) and
\(\lVert  \widetilde{g}_i  \rVert_\infty \le c \coloneqq 36 \, C_{\mathcal{Q}}^2 \, \lVert  u  \rVert^2 \, t_{\text{mix}}^2 \, \lVert  \varepsilon  \rVert_\infty^2\).
Apply \eqref{eq:external-markov-conc} with \(c_i = c\):
\begin{equation}
\begin{aligned}
\mathbb{E}_\xi^{1 / p} \left[
  \left(|\sum_{l = 2}^{n - 1} g_l(Z_{l - 1})|\right)^p
\right]
  \le C_{\text{MC}} \sqrt{p \, t_{\text{mix}} \, (n - 2) c^2}.
\end{aligned}
\end{equation}
Substitution gives
\begin{equation}
\begin{aligned}
\mathbb{E}_\xi^{1 / p} \left[
  \left(|\sum_{l = 2}^{n - 1} g_l(Z_{l - 1})|\right)^p
\right]
  \le C_{\text{MC}} \cdot 36 \, C_{\mathcal{Q}}^2 \, \lVert  u  \rVert^2
       \, \lVert  \varepsilon  \rVert_\infty^2 \, t_{\text{mix}}^{5 / 2}
       \, \sqrt{p \, n}.
\end{aligned}
\end{equation}
Absorb the universal prefactor into \(C_4\). \(\square\)

\begin{corollary}\label{cor:M-RR-bracket-asymp}
  \textbf{(Stationary asymptotic-variance bracket comparison.)}
  Under the assumptions of the previous lemma, for every \(u \in \mathbb{R}^d\), every
  \(p \ge 2\), and every \(n \ge 2\),
\begin{equation}
\begin{aligned}
  \mathbb{E}_\xi^{1 / p} \left[
    | u^\top \langle M^{\text{RR}} \rangle_n u - n \, \sigma^2(u) |^p
  \right]
    \le C_4 \, C_{\mathcal{Q}}^2 \, \lVert  u  \rVert^2 \, \lVert  \varepsilon  \rVert_\infty^2
       \, t_{\text{mix}}^{5 / 2} \, \sqrt{p \, n}
       + \frac{C_3 \, \lVert  u  \rVert^2}{\alpha \, a},
\end{aligned}
\label{eq:M-RR-conc-sigma}
\end{equation}
  with \(C_3\) the variance-comparison constant of Section 4.5.
\end{corollary}

% \textit{Proof.} Use the triangle inequality:
% \begin{equation}
% \begin{aligned}
% | u^\top \langle M^{\text{RR}} \rangle_n u - n \, \sigma^2(u) |
%   \le | u^\top \langle M^{\text{RR}} \rangle_n u - n \, \sigma_n^{2, \text{RR}}(u) |
%    + n \, | \sigma_n^{2, \text{RR}}(u) - \sigma^2(u) |.
% \end{aligned}
% \end{equation}
% The first piece is bounded by \eqref{eq:M-RR-conc}; the second by
% Lemma~\ref{lem:RR-variance-comparison}.
% \(\square\)
